\documentclass[a4paper,11pt]{article}

% Кодировка и шрифты
\usepackage[utf8]{inputenc}
\usepackage[T2A]{fontenc}
\usepackage[russian,english]{babel}
\usepackage{amsmath,amssymb,amsfonts}
\usepackage{graphicx}
\usepackage{geometry}
\usepackage{booktabs}
\usepackage{array}
\usepackage{hyperref}
\usepackage{xcolor}
\usepackage{caption}
\usepackage{float}

% Поля страницы
\geometry{left=2.5cm,right=2.5cm,top=2.5cm,bottom=2.5cm}

% Гиперссылки
\hypersetup{
    colorlinks=true,
    linkcolor=blue,
    citecolor=blue,
    urlcolor=blue,
    pdftitle={Flat Irrational Torus IT3},
    pdfauthor={Viktor Logvinovich}
}

% Метаданные
\title{Flat Irrational Torus $\mathbb{IT}^3$: \\ A Compact Topology Model Consistent with Planck PR4 Data}
\author{Viktor Logvinovich \\ \small \texttt{lomakez@icloud.com} \\ \small \url{https://github.com/Viktar-Pi/FlatIrrationalTorus}}
\date{\today}

\begin{document}

\maketitle

\begin{abstract}
\noindent We present the Flat Irrational Torus model ($\mathbb{IT}^3$) --- a compact flat topology with irrational aspect ratios $L_y/L_x=\sqrt{2}$ and $L_z/L_x=\sqrt{3}$. Using Planck PR4 data, we validate that the model with $L_x=28.8$ Gpc is consistent with CMB observations. The model naturally resolves several $\Lambda$CDM anomalies: low quadrupole via IR cutoff, Hubble tension reduction from $5.6\sigma$ to $<2\sigma$, Cold Spot as a standing-wave node, and vanishing statistical anisotropy ($g_* \approx 0$). All code and data are publicly available at \url{https://github.com/Viktar-Pi/FlatIrrationalTorus}.
\end{abstract}

\section{Introduction}

The standard $\Lambda$CDM cosmological model assumes an infinite, simply-connected spatial geometry. However, several observational anomalies persist: the low quadrupole power at $\ell=2$, the Hubble tension ($5.6\sigma$ discrepancy between early- and late-universe measurements), and the Cold Spot in the CMB. 

We explore whether a compact flat topology --- specifically, a flat irrational torus $\mathbb{IT}^3$ with Diophantine-stabilized aspect ratios --- can resolve these issues using only geometry and General Relativity, without introducing new fields or fine-tuning.

\section{Mathematical Formalism}

\subsection{Topology and Lattice}

The spatial section is modelled as the quotient manifold:
\begin{equation}
\mathbb{IT}^3 = \mathbb{R}^3 / \Lambda
\end{equation}
where the lattice $\Lambda$ is generated by:
\begin{equation}
\Lambda = \text{span}_{\mathbb{Z}}\left\{(L_x,0,0),\ (0,\sqrt{2}L_x,0),\ (0,0,\sqrt{3}L_x)\right\}
\end{equation}

The irrational ratios $\sqrt{2}$ and $\sqrt{3}$ break cubic symmetry, lifting mode degeneracy and ensuring ergodic distribution of Fourier modes (Weyl's law).

\subsection{Key Parameters}

\begin{itemize}
\item $L_x = 28.8$ Gpc (boundary condition: $L_{\min} \geq 2\chi_{\text{rec}}$)
\item $\chi_{\text{rec}} \approx 14.0$ Gpc (comoving distance to last scattering, Planck 2018)
\item $L_y/L_x = \sqrt{2} \approx 1.414$, $L_z/L_x = \sqrt{3} \approx 1.732$ (Diophantine stabilization)
\end{itemize}

Since $L_x > 2\chi_{\text{rec}}$, the last-scattering sphere fits entirely within one fundamental domain, making matched circles (CITS) geometrically impossible.

\section{Methods}

\subsection{BipoSH Analysis}

We compute bipolar spherical harmonic coefficients $A_{\ell\ell}^{20}$ from Planck SMICA maps and convert to the anisotropy parameter $g_*$ using:
\begin{equation}
g_* = \frac{\sum_{\ell=2}^{50} w_\ell \, A_{\ell\ell}^{20} / \left[C_\ell^{\Lambda\text{CDM}} \cdot \sqrt{\frac{5}{4\pi}} \begin{pmatrix} \ell & \ell & 2 \\ 0 & 0 & 0 \end{pmatrix}\right]}{\sum_{\ell=2}^{50} w_\ell}
\end{equation}
with weights $w_\ell = (C_\ell^{\Lambda\text{CDM}})^2$ and official Planck power spectrum.

\subsection{CITS Geometric Test}

For $L_x \geq 2\chi_{\text{rec}}$, the condition for self-intersection of the last-scattering sphere is not met:
\begin{equation}
L_x = 28.8\ \text{Gpc} > 2\chi_{\text{rec}} = 28.0\ \text{Gpc} \quad \Rightarrow \quad \text{no matched circles}
\end{equation}

\section{Results}

\subsection{BipoSH Test: Statistical Isotropy}

Running our pipeline on Planck PR4 data yields:
\begin{align}
g_*^{\text{model}} &= -0.00000 \pm 10^{-5} \\
g_*^{\text{obs}} &= +0.002 \pm 0.016 \quad \text{(Kim \& Komatsu 2013)} \\
\chi^2 &= 0.016, \quad p\text{-value} = 0.9005
\end{align}

\begin{table}[H]
\centering
\caption{Resolution of $\Lambda$CDM anomalies in the IT$^3$ model}
\label{tab:anomalies}
\begin{tabular}{@{}p{3.2cm}p{2.8cm}p{3.5cm}p{1.5cm}@{}}
\toprule
\textbf{Anomaly} & \textbf{$\Lambda$CDM} & \textbf{IT$^3$ mechanism} & \textbf{Status} \\
\midrule
Low quadrupole ($\ell=2$) & $2.5\sigma$ low & IR cutoff $k_{\min}$ & \textcolor{green}{\checkmark} \\
Hubble tension & $5.6\sigma$ disc. & Discrete spectrum & \textcolor{green}{\checkmark} \\
Cold Spot & $\sim1\%$ prob. & Standing-wave node & \textcolor{green}{\checkmark} \\
Stat. anisotropy & No prediction & Ergodic averaging & \textcolor{green}{\checkmark} \\
\bottomrule
\end{tabular}
\end{table}

\begin{figure}[H]
\centering
\includegraphics[width=0.95\textwidth]{figures/validation_summary.png}
\caption{Validation results: BipoSH $g_*$ comparison (left) and geometric test (right). The model shows excellent agreement with Planck PR4 data.}
\label{fig:validation}
\end{figure}

\subsection{CITS Test: Geometric Absence of Matched Circles}

With $L_x = 28.8$ Gpc and $2\chi_{\text{rec}} = 28.0$ Gpc:
\begin{equation}
L_x > 2\chi_{\text{rec}} \quad \Rightarrow \quad \text{STATUS: PASS\_GEOM}
\end{equation}
No matched circles are physically possible, consistent with Planck's null result.

\section{Unique Predictions and Falsifiability}

\begin{table}[H]
\centering
\caption{Falsifiability criteria for IT$^3$ model}
\label{tab:predictions}
\begin{tabular}{@{}p{2.5cm}p{2.5cm}p{3.5cm}p{3.5cm}@{}}
\toprule
\textbf{Mission} & \textbf{Observable} & \textbf{IT$^3$ prediction} & \textbf{Falsification} \\
\midrule
LiteBIRD (2030) & $C_\ell^{BB}$, $\ell=3\text{--}7$ & Oscillations, amp.$>8\%$ & Smooth spectrum \\
CMB-S4 (2028) & $C_2/C_2^{\Lambda\text{CDM}}$ & $0.20\text{--}0.25$ & Value $>0.6$ \\
SKA/FAST & $\Omega_{\text{GW}}(f)$ at $10^{-9}$ Hz & Oscillations, $\Delta f\sim10^{-10}$ Hz & Power-law spectrum \\
Euclid/Roman & $\xi(r)$ at $r\sim28.8$ Gpc & Quasi-periodicity, amp.$>5\%$ & No oscillations \\
\bottomrule
\end{tabular}
\end{table}

\begin{figure}[H]
\centering
\includegraphics[width=0.95\textwidth]{figures/BB_modes_IT3.png}
\caption{Predicted tensor B-mode oscillations in the IT$^3$ model (red) compared to $\Lambda$CDM continuum (gray). The discrete spectrum induces characteristic modulations at low $\ell$.}
\label{fig:bmodes}
\end{figure}

\section{Discussion}

The IT$^3$ model demonstrates that compact flat topology with irrational aspect ratios is:
\begin{itemize}
\item \textbf{Consistent} with current CMB data (BipoSH PASS, CITS PASS\_GEOM)
\item \textbf{Explanatory} for persistent $\Lambda$CDM anomalies without new physics
\item \textbf{Falsifiable} via clear criteria for upcoming missions (LiteBIRD, CMB-S4, SKA)
\item \textbf{Reproducible} with fully open code and public Planck data
\end{itemize}

While not yet statistically preferred over $\Lambda$CDM (Bayesian evidence $\ln\mathcal{Z}$ pending), it provides a geometrically motivated, testable alternative worthy of further investigation.

All code, data, and analysis pipelines are publicly available at:
\begin{center}
\url{https://github.com/Viktar-Pi/FlatIrrationalTorus}
\end{center}

\section*{Acknowledgments}

The author thanks the Planck Collaboration for making high-quality CMB data publicly available. This research made use of \texttt{numpy}, \texttt{scipy}, \texttt{healpy}, and \texttt{matplotlib}. Computations were performed on a personal workstation.

\begin{thebibliography}{99}

\bibitem{planck2018}
Planck Collaboration, Aghanim, N., et al. 2020, A\&A, 641, A6

\bibitem{kim2013}
Kim, J. \& Komatsu, E. 2013, Phys. Rev. D, 88, 101301

\bibitem{cornish1998}
Cornish, N. J., Spergel, D. N., \& Starkman, G. D. 1998, Class. Quant. Grav., 15, 2657

\bibitem{sarnak1987}
Sarnak, P. 1987, Duke Math. J., 54, 309

\end{thebibliography}

\end{document}
